% =============================================================================
% INTRODUCTION TO VOLUME 8: THE STANDARD MODEL FROM GF(27), LEPTON MASSES AND
% INTERFACES (Vol. 8 of 8)
% =============================================================================

\section*{Eighth Volume of the Document Collection}

Volume~8 comprises Documents 344 to 389 from September to early October 2026
and brings the collection up to the current state of the corpus. It builds
directly on the Galois bridge of Volume~7 and carries it forward to the
quantum numbers of the Standard Model, to the lepton masses, and to the
interfaces with other theoretical frameworks.

\subsection*{Content and Focus}

The first part derives from the structure of $GF(27)^*$ the quantisation of
charge, the Gell-Mann--Nishijima relation $Q=I_3+Y/2$, the generation
structure with its coupling matrix, and $SU(2)_L$ chirality; this is
complemented by the comparison with Furey fermions, a placement within the
literature, and the treatment of black holes under time--mass duality. The
second part first clarifies how the PDG comparison values arise and which
model assumptions enter them, and then sets the lepton ladder, the Galois
approach and the Koide relation side by side; there follow the Koide
amplitude $d/c=\sqrt2$, the Higgs geometry, spontaneous symmetry breaking from
the $T^4/\mathbb{Z}_3$ geometry, and the overview of the Standard Model
particles with their certificate status.

The third part shows why harmonics and algebraic structure are the same
mathematics, and tests this insight on laboratory systems: frequency analysis
and its limits, HRV measurement, the kagome lattice and the BAW Ising machine.
Added to these are Hodge theory on $T^4/\mathbb{Z}_3$, the comparison of the
icosahedral carrier with the $D_4$ lattice, the presentation of the foundation
in four layers, and the field structure of energy, flux and geometry. The
fourth part is devoted to the lepton masses: $\theta=p_0=2/9$ as two readings
of one quotient, the route from $p_0,p_1,p_2$ to the masses, four routes
compared, the justification of the icosahedron, the factor $4/3$, the
matrix-element principle, and the Yukawa mechanism as a consequence of
$\tilde T\cdot m=1$.

The fifth part collects the interfaces: the bridge equations to Observer Patch
Holography, the hierarchy $v/E_P$, an overview of all comparison and bridge
documents, considerations on peer review, the specificity of the $D_4$
carrier, the exactly summed recursion, the status of muon $g-2$, the masses
without $v$, a brief account of FFGFT after the calculation review, the
Higgs--vacuum route to $\xi$, the place of $\alpha$ in the geometry, an
overview of the deviations, the CMB temperature in the basic form, and the
representations of gravitation using the example of parabolic flight.

\subsection*{Organisation}

The 46 documents are divided into five parts: quantum numbers from the Galois
structure (Docs.~344--350); comparison values, lepton residuals, Higgs and
Standard Model particles (Docs.~351--357); harmonics, laboratory systems,
Hodge theory and fields (Docs.~358--366); icosahedron, matrix elements and
lepton masses (Docs.~367--373); interfaces, review and the basic form
(Docs.~374--389).

\subsection*{Reading the Documents}

For a quick overview of the overall state, Doc.~384 (brief account) is
suitable; for the remaining deviations, Doc.~387. Status markers and the
corrections register apply as in the preceding volumes; in each case the most
recent version of a result in the corpus is authoritative.

\subsection*{Status}

The documents in this volume reflect their respective versions in the corpus
as of October 2026.
